\begin{table}[htbp]
\centering
\small
\caption{Escolha da melhor métrica few-shot do Problema A para a transferência (openai/gpt-4o-mini).}
\label{tab:T7_escolha_gpt4mini}
\begin{tabular}{lrrrrrrrl}
\toprule
Métrica de origem (A) & Fid.\ no A & Métr.\ vs real (A) & cos(Perc., NNLS) & Consist.\ B & $\kappa$ B & Consist.\ C & $\kappa$ C & Eleição \\
\midrule
zero-shot (referência) & 81,3 (78,0–85,3) & 34,7 (17,3–49,3) & 0,980 & 70,7 & 0,348 & 88,0 & 0,745 &  \\
4-shot & 89,2 (84,2–94,5) & 86,1 (68,0–94,0) & 0,990 & 70,7 & 0,354 & 88,7 & 0,760 &  \\
10-shot & 94,4 (90,7–98,6) & 93,4 (86,7–98,0) & 0,833 & 73,2 & 0,407 & 87,9 & 0,745 &  \\
\textbf{20-shot} & \textbf{95,5 (92,3–97,7)} & \textbf{93,9 (89,3–98,0)} & \textbf{0,789} & \textbf{74,9} & \textbf{0,440} & \textbf{87,6} & \textbf{0,739} & \textbf{eleita} \\
40-shot & 95,2 (90,9–98,2) & 93,4 (88,0–98,0) & 0,939 & 77,1 & 0,487 & 87,3 & 0,735 &  \\
\bottomrule
\multicolumn{9}{p{0.95\linewidth}}{\footnotesize Critério de eleição: maior fidelidade média no Problema A (Perceptron, R$^2$, sem colapso) entre as métricas few-shot; o zero-shot entra só como referência e o teto não concorre. Consistência e $\kappa$ em B e C: rótulos zero-shot do LLM contra a métrica de origem (Perceptron, só $\hat{W}$ transferido, centróides recalculados pelos rótulos do LLM). Em negrito: a métrica eleita para a transferência.} \\
\end{tabular}
\end{table}
